% ============================================================================
% RELATED_WORK_DRAFT.tex — rewritten Related Work (REVIEW COPY, data-independent).
% Flowing four-movement narrative grounded in the verified LIT_SURVEY.md. Each
% movement closes on the precise gap our shared-camera coupling occupies. Reuses the
% bib keys already in ACODS.tex where they match; NEW keys needed are listed at the
% bottom (all verified in LIT_SURVEY.md). Framing-insensitive: leads with the field,
% not with our contribution ordering, so it stands regardless of the deferred
% lead-identity decision.
% ============================================================================

\section{Related Work}
\label{sec:related}

\paragraph*{Perception-aware and active-perception control.}
A growing body of work rejects the classical separation of perception from control,
observing that a robot's motion determines what it can sense and should therefore be
chosen to serve estimation, not only the task. Perception-aware model-predictive
control couples visibility of a point of interest to the quadrotor dynamics through the
cost~\cite{falanga2018pampc}, differential-flatness formulations keep a target in the
field of view during aggressive flight~\cite{spasojevic2020perception}, and active
visual localization shapes motion to minimize pose uncertainty, e.g.\ by querying a
precomputed Fisher-information field~\cite{zhang2019fisher}. This line inherits the
active-perception thesis that sensing, motion, and estimation form a single
loop~\cite{cadena2016past,placed2023survey}. In every case, however, the perceiving
agent and the acting agent are the \emph{same} robot, perception quality is
\emph{monotone} in the chosen viewpoint---a better viewpoint simply lowers the robot's
own covariance---and the decision variable is the ego-robot's trajectory or gaze. We
retain the co-optimization principle but change its structure: a \emph{single shared
sensor} carried by one agent supplies both the navigation reference and the state
estimate of a \emph{second} agent, and a single scalar---altitude---moves the two
perception products in \emph{opposite} directions.

\paragraph*{Belief-space planning under estimation uncertainty.}
Our UGV controller belongs to the belief-space family, which augments the planning
state with the estimator covariance so that motion is planned under partial
observability. Representative formulations factor the covariance for efficient
roadmap planning~\cite{prentice2009belief}, optimize locally in Gaussian belief
space~\cite{vandenberg2012motion}, and exploit maximum-likelihood observations for
tractability~\cite{platt2010belief}, with generalized-belief-space extensions to
unknown environments~\cite{indelman2015planning} and roots in planning under
uncertainty~\cite{kaelbling1998planning}. The defining move of this literature is that
the robot shapes its \emph{own} future observations---the canonical ``move into the
light'' detour to localize before committing to a goal. Our setting differs
structurally: the informativeness of the UGV's observations is not set by its own
motion but by an \emph{external teammate's} altitude, which the UGV does not control,
and that same altitude simultaneously fixes the navigation reference the planner must
track. We thus inherit the belief-space machinery but relocate the information-design
decision from the ego-robot's trajectory to a cooperating agent's sensor geometry.

\paragraph*{Cooperative aerial--ground teams.}
Using an aerial vehicle as an ``eye in the sky'' for a perception-limited ground robot
is well established~\cite{grocholsky2006cooperative}, with recent systems guiding a
perception-limited~\cite{li2024colag} or perception-less~\cite{santos2024aerial} ground
vehicle from the air, alongside quadrotor target tracking under field-of-view and
occlusion constraints~\cite{thomas2017autonomous,penin2018vision} and EKF-based
cooperative localization~\cite{roumeliotis2002distributed,kia2016cooperative}. A
recurring theme is a coverage--accuracy synergy---the aerial platform provides wide
coverage while the ground platform provides accurate, close-range
sensing~\cite{grocholsky2006cooperative}. Crucially, that trade-off is resolved by
assigning coverage and accuracy to \emph{different} robots or \emph{different} sensors,
and the aerial altitude is fixed and treated as a passive relay parameter. We show
instead that when a \emph{single} downward camera serves both roles, altitude becomes
an active perception control input whose optimal value is dictated---in closed form and
as a function of path curvature (Prop.~\ref{prop:alt}, Cor.~\ref{cor:climb})---by the
shared-camera coupling itself.

\paragraph*{Event-triggered control and inter-execution guarantees.}
Event- and self-triggered control re-solve only when a state-dependent condition
fires, curbing the computation and communication of periodic
execution~\cite{tabuada2007event,heemels2012introduction}, with dynamic triggering that
enlarges inter-event intervals~\cite{girard2015dynamic} and extensions to robust MPC of
nonlinear systems~\cite{eqtami2010event,li2014event}. Boundedness of the inter-execution
error rests on input-to-state stability; the discrete-time theory we invoke is due to
Jiang and Wang~\cite{jiangwang2001}, and evaluating the trigger on a periodic clock
yields a minimum inter-event time by construction, precluding Zeno
behavior~\cite{heemelsdonkersteel2013}. Triggers in this literature are almost always
functions of the \emph{state} error; the closest precedent to a perception-driven
trigger is variance-based event-triggered \emph{estimation}, which fires when a
predicted error covariance exceeds a threshold~\cite{trimpe2014event}. Our trigger
shares the idea of firing on an estimation quantity but differs in three respects: it
is a \emph{geometric} function of altitude (footprint/GSD), it fires a control
\emph{re-plan} rather than an estimator transmission, and it operates over a
disturbance whose covariance $R(z)\propto z^{2}$ is itself \emph{set by the altitude
being chosen}---so the trigger is one half of a control-and-altitude co-design rather
than a fixed-plant monitor.

\paragraph*{Robust tube-based control.}
The ancillary tracker draws on tube MPC, which splits control into a nominal plan and
an ancillary law confining the true state to a robust positively invariant
tube~\cite{mayne2005robust,langson2004tubes,chisci2001persistent}; the standard outer
approximation of the minimal invariant set~\cite{mRPI} and the nonlinear
case~\cite{mayne2011tube} are computed offline for a fixed disturbance set. We use this
design only to fix the ancillary gain and an ellipsoidal shape matrix; the online stage
replaces the polyhedral tube with a lightweight adaptive ellipsoidal surrogate
(Sec.~\ref{sec:method}), trading formal recursive feasibility for a real-time
perception-driven monitor whose inter-trigger boundedness is nonetheless certified by
Prop.~\ref{prop:iss}.

\medskip\noindent\emph{Gap.} Across these literatures, perception is optimized by a
single robot for its own estimate, belief-space information is gathered by the
ego-robot's own motion, aerial--ground coverage and accuracy are split across separate
robots or sensors, and event triggers monitor a fixed plant. None treats a single
shared camera's altitude as one control variable that antagonistically sets both
localization resolution and navigation coverage, with a closed-form curvature-scheduled
optimum realized by a perception-triggered altitude/gain co-design carrying
discrete-time boundedness and no-Zeno guarantees. That is the space this paper occupies.

% ---------------------------------------------------------------------------
% NEW bib keys required (verified in LIT_SURVEY.md); add to ICC_kisan.bib:
%   girard2015dynamic       A. Girard, "Dynamic triggering mechanisms for
%       event-triggered control," IEEE TAC 60(7):1992-1997, 2015.
%   jiangwang2001           Z.-P. Jiang, Y. Wang, "Input-to-state stability for
%       discrete-time nonlinear systems," Automatica 37(6):857-869, 2001.
%   heemelsdonkersteel2013  W.P.M.H. Heemels, M.C.F. Donkers, A.R. Teel, "Periodic
%       event-triggered control for linear systems," IEEE TAC 58(4):847-861, 2013.
%   trimpe2014event         S. Trimpe, R. D'Andrea, "Event-based state estimation with
%       variance-based triggering," IEEE TAC 59(12):3266-3281, 2014.
%   ranvreugdenhil1988      A.C.M. Ran, R. Vreugdenhil, "Existence and comparison
%       theorems for algebraic Riccati equations...," Lin. Alg. Appl. 99:63-83, 1988.
%   andersonmoore1979       B.D.O. Anderson, J.B. Moore, "Optimal Filtering,"
%       Prentice-Hall, 1979.
% Existing ACODS.tex keys reused above (verify they resolve in ICC_kisan.bib):
%   falanga2018pampc, spasojevic2020perception, zhang2019fisher, cadena2016past,
%   placed2023survey, prentice2009belief, vandenberg2012motion, platt2010belief,
%   indelman2015planning, kaelbling1998planning, grocholsky2006cooperative,
%   li2024colag, santos2024aerial, thomas2017autonomous, penin2018vision,
%   roumeliotis2002distributed, kia2016cooperative, tabuada2007event,
%   heemels2012introduction, eqtami2010event, li2014event, mayne2005robust,
%   langson2004tubes, chisci2001persistent, mRPI, mayne2011tube
% ============================================================================
